%date --- 12 - 08 -26

\subsection{Introduction}

	We start with the assumption that you have no idea what General Relativity is.
	How do we quantify what is gravity? What even is gravity? We begin with the classical equations of motion.

	To motivate ourselves, we describe three concepts introduced over the history of classical mechanics and explain how 
	it leads to certain contradictions or problems that motivated the invention of general relativity

	\subsubsection{Newton's Law of Gravitation}

		Newton's Laws as he prescribed them is expressed as follows:

		\begin{equation}
			\vb{F} = -\frac{GMm}{r^2}\vb{\hat{r}} 
		\end{equation}

		This is the familiar form of Newtonian gravity, and is perhaps one of the most recognizable equations in physics.
		This is what every high schooler knows and what every mentor presents, and is trivial.

		We can separate the mass $m$ experiencing the force from the mass $M$ producing it, and write

		\begin{equation}
			\vb{F} = m_g \qty(-\frac{GM}{r^2}\hat{\vb{r}})
		\end{equation}

		We can henceforth call $m_g$ as the \textit{gravitational mass} and the quantity $-\frac{GM}{r^2}\hat{\vb{r}}$ as the \textit{gravitational field}

		We can express it more generally by recognizing the force field is simply the negative gradient of a potential field.
		\begin{equation}
			\vb{F} = m_g\qty(-\grad\Phi)
			\label{eq: F = mg x -gradphi}
		\end{equation}
		where $\displaystyle \Phi = -\frac{GM}{r} $ is the gravitational potential produced by the source mass.

		Thus, we can henceforth call $m_g$ as the \textit{gravitational charge}. It follows from the fact that
		\begin{equation}
			\text{charge}\times \text{field gradient} = force
		\end{equation}
		This formulation is universal across all of physics.

		Analogous to electromagnetism, where 
		\begin{equation}
			\vb{F} = \frac{1}{4\pi\varepsilon_0} \frac{q_1 q_2}{r^2}\hat{\vb{r}} 
		\end{equation} 
		Like the electric charge $q$, the gravitational charge $m_g$ determines the strength of the interaction of the forces.

		An important property that Newtonian gravitation has is that it has superposition. i.e., fields add together linearly.
		This is therefore also true of the gravitational potential.

		For a system of point masses, the total potential field $\Phi(\vb{x})$ at a point $\vb{x}$ is given by the sum of the potentials produced by each individual mass:

		\begin{align*}
			\Phi(\vb{x}) &= \sum_i \Phi_i \\
							&= -G\sum_i \frac{M_i}{\abs{\vb{x}-\vb{r}_i}} 
		\end{align*}

		We can extend this description from a discrete collection of point masses to a continuous mass distribution.
		Let the mass contained in an infinitesimal volume $\dd[3]{x}$ be

		\begin{equation}
			dm_i = \rho(t,\vb{x}) \dd[3]{x}
		\end{equation}
		Then we arrive with the continuum version of the gravitational potential field generated by a mass distribution.
		\begin{equation}
			\Phi(\vb{x}) = - G \int \frac{\rho(t,\vb{x})\dd[3]{r}}{\abs{\vb{x}-\vb{r}}}
		\end{equation}

		Here, $\vb{r}$ denotes the position of each source element of the mass distribution, while $\vb{x}$ denotes the point at which the potential is being evaluated.

		We can now obtain a differential equation for the gravitational potential. Taking the Laplacian of $\Phi$ gives

		\begin{equation}
			\lpc \Phi = -G \int \rho(t,\vb{x})\dd[3]{r} \lpc \frac{1}{\abs{\vb{x}-\vb{r}}}
		\end{equation}
		Recall that this operation results to a dirac-delta function
		\begin{equation}
			\lpc \frac{1}{\abs{\vb{x}-\vb{r}}} = - 4\pi\delta(\vb{x}-\vb{r})
		\end{equation}

		which relates the Laplacian of the inverse-distance potential to the Dirac delta function. Substituting this into the previous expression gives

		\begin{equation}
			\lpc \Phi = 4\pi G \int \rho(t,\vb{x})\dd[3]{r} \delta(\vb{x}-\vb{r})
		\end{equation}
		Recall the sifting property of the Dirac delta
		\begin{equation}
			\int f(x)\delta(x-x') \dd{x'} = f(x')
		\end{equation}
		From this we obtain	
		\begin{equation}
			\lpc \Phi = 4\pi G \rho(t,\vb{r})
		\end{equation}
		Thus, we arrived with the fundamental field equation which governs Newtonian Gravity

		\begin{definition}{Field Equation for Newtonian Gravity}
			\begin{equation}
			   \lpc \Phi = 4\pi G \rho  
			\end{equation}
		\end{definition}
		The significance of a field equation is that it relates a field to its sources. In this case, the mass density $\rho$ acts as the source of the gravitational potential field $\Phi$. 
		
		Note the following
		\begin{equation}
			\Phi = -G\int \frac{\rho \dd[3]{r}}{\abs{\vb{x}-\vb{r}}} 
			\iff \lpc\Phi = 4\pi G \rho 
		\end{equation}

		The two are notionally equivalent. One is a differential equation and the other is the integral form. One can switch
		to and from them using the Laplacian and Green's functions\footnote{Green's functions are essentially the inverse of a laplacian, $\nabla^{-2}$}.
		It assigns the definition of the field, and fields assign values to every point in space.

		However, this formulation presents a fundamental problem. Since the potential $\Phi$ is determined directly by the mass density $\rho$, any change in the mass distribution produces 
		an immediate change in the gravitational field throughout space. There is no propagation delay or retardation. In other words, Newtonian gravity permits changes in the gravitational
		field to propagate instantaneously, which is incompatible with relativistic causality. This problem is one of the motivations for general relativity. 

		The corresponding requirement in GR is that, in the appropriate weak-field and slow-motion limit, general relativity must reduce to Newtonian gravity.

		Now that we have the potential of the field that the source creates, and thus the force, how do we determine how an object moves in response to that field?

	\subsubsection{Newton's Equations of Motion}
		The quintessential equation governing the motion of classical bodies was expressed by Newton as
		\begin{equation}
			\vb{F} = m\vb{a}
		\end{equation}
		Everybody and their mother knows this shit.

		This states that if one sums of all the forces acting on an object, 
		it must equal its mass multiplied by its acceleration:
		\begin{equation}
			\sum\vb{F} = m\vb{a}
		\end{equation}
		We can rewrite the equation as
		\begin{equation}
			\vb{a} = \frac{1}{m} \sum \vb{F}
			\label{eq: a = 1/m f}
		\end{equation}

		Notice that the $m$ term acts as a kind of resistance term\footnote{Recall Ohm's Law $V = IR$.}, similar to the concept of inertia. This is to be expected, since Newton's first law can essentially be understood as the statement that, in the absence of a net force, an object does not change its state of motion.

		Hence, we name this quantity the \textit{inertial mass}, $m_i$. Inertia denotes the difficulty or resistance of an object to a change in its state of motion.

		We can restate (\ref{eq: a = 1/m f}) with $m$ replaced by $m_i$, to distinguish inertial mass from the gravitational mass introduced previously:
		\begin{equation}
			\vb{a} = \frac{1}{m_i} \vb{{F}}
		\end{equation}

		Recall that we already have the definition of $\vb{F}$ from (\ref{eq: F = mg x -gradphi}). Thus, to calculate the motion of an object under the influence of gravity, we simply calculate its acceleration:
		\begin{equation}
			\vb{a} = \frac{m_g}{m_i}\qty(-\grad\Phi)
		\end{equation}

		Notice that this is similar in form to the equation of motion for a charged particle in an electric field.
		\begin{equation}
			\vb{a} = \frac{q}{m_i}\vb{E}
		\end{equation}
		where $m_i$ is the inertial mass, the measure of an object's resistance to changes in its motion. In both cases, the coupling to the field determines the force, while the inertial mass determines the resulting acceleration.

		However, for calculations involving Newtonian gravity, the ratio of the gravitational charge and inertial mass
		$\displaystyle \frac{m_g}{m_i}$ is assumed to be $1$. Thus, for gravity, the acceleration of a body is expressed as
		\begin{equation}
			\vb{a} = -\grad\Phi.
		\end{equation}
		This assumes that
		\begin{equation}
			m_g = m_i,
		\end{equation}
		This feels unnatural and unexpected, a truly remarkable coincidence in Newtonian theory. Why should it be the case that one force acts differently from other forces in terms of its charge carriers?
		Mass is a fundamental property that affects other forces, there's no immediate evidence why mass is gravity's charge carrier.

		There is no \textit{a priori} reason why the magnitude of gravitational force on a particle should equal the quantity that determines the particle's resistance to an applied force in general. 
		They play fundamentally different roles in the theory. It is 
		analogous to assuming that an object's coupling to some other field must be equal to 
		its inertial mass, like an electron's electromagnetic properties being dependent on mass.

		Nevertheless, for gravity this equality is accepted experimental fact without explanation. 
		\newpage

	\subsubsection{Weak Equivalence Principle}
		The assumption that
		\begin{equation}
			m_i=m_g
		\end{equation}
		is called the \textit{weak equivalence principle}.

		It follows that, for objects \underline{only} affected by gravity, the equation of motion is simply
		\begin{equation}
			\vb{a} = -\grad\Phi.
		\end{equation}

		That is, any test mass placed in the same gravitational field experiences the same acceleration, regardless of its mass or composition. The nature of the object itself does not appear in the equation of motion. The field prescribes the acceleration, the test body merely responds to it.

		One could therefore remove the test object $m$ and nothing about the gravitational field would change. In Newtonian mechanics, this is treated as a universal fact: the gravitational field exists independently of the particular body used to probe it.

		However, this raises a deeper question.

		One could argue that the stage for the motion of an object is the field itself. If we remove the object from some local region, the field remains.

		By the same logic, if we remove the field itself, what remains? In Newtonian mechanics, space $\R^3$ remains as the stage on which the field and matter exist.

		The influence of gravity has \underline{nothing} to do with the nature of the test object itself. The gravitational field does not change when the test object is removed, and its motion depends only on the field at its location.

		One might ask: if we remove all mass, would there not be no gravity? The answer is that we must distinguish the sources of the field from the test body. We are considering local interactions. Mass may exist somewhere outside the local region of interest and still generate a gravitational field within it. Given the local potential $\Phi$, we can determine information about its sources, but the test body itself is not one of those sources.

		Thus,
		\begin{equation}
			\vb{a} = -\grad\Phi
		\end{equation}
		is treated as inertial motion in a gravitational field.

		In the stage analogy, one could think of $-\grad\Phi$ as living on the stage of space $\R^3$. The field exists on space and determines how objects move through it. 
		In some fields such as particle physics, that is the case, they treat gravity as a special field  which acts on every particle and every particle responds to in the same way

		However, in GR, $-\grad\Phi$ is no longer treated as a field sitting on top of the stage. The gravitational field \textit{is} the stage. 
		The space time is what the rest of physics takes place in.

		The strong equivalence principle extends this idea: locally, the laws of physics in a freely falling frame are the same as those of an inertial frame in the absence of gravity. In other words, locally, gravity can be transformed away by choosing an appropriate freely falling frame.

		Inertial motion is therefore determined by the structure of spacetime.

		Physics takes place on the stage called spacetime, and GR governs that space time. The spacetime structure remains even after the particular body used to probe it is removed.

		In GR, instead of the gravitational potential field $\Phi$, we have the metric:
		\begin{equation}
			ds^2 = g_{\mu\nu} dx^\mu dx^\nu.
		\end{equation}

		The metric determines the geometry of spacetime and, consequently, the inertial trajectories of objects.

		In particle physics, one generally begins with a fixed flat spacetime $\R^4$ and defines fields upon that background. Gravity can then be treated as a field on this background, rather than as the geometry of spacetime itself. This approach is sound as a field-theoretic formulation and provides a natural route toward quantizing gravity and describing gravitons, but it is ugly and inelegant. Regardless, you can bootstrap yourself from there and arrive on the EFE

		The distinction of perspectives can be summarized as follows:

		\begin{center}
			\begin{tikzpicture}[>=stealth, node distance=2.5cm]
				\node (rel) {
					\begin{tabular}{c}
						\textbf{Relativist perspective}\\[2mm]
						Spacetime structure is variant\\
						$\downarrow$\\
						Gravity is the manifestation\\ 
						of Riemannian space time\\ 
						curvature and it governs\\
						particle motion
					\end{tabular}
				};

				\node (part) [right=4cm of rel] {
					\begin{tabular}{c}
						\textbf{Particle physicist perspective}\\[2mm]
						Fixed spacetime, \\fields live on top of it\\
						$\downarrow$\\
						Gravity is a special \\gauge field in $SO(3,1)$\\
						which couples to\\ everything universally
					\end{tabular}
				};

				\node (metric) [above=1.2cm of $(rel)!0.5!(part)$] {
					\textbf{Metric}
				};

				\draw[->] (metric) -- (rel);
				\draw[->] (metric) -- (part);
			\end{tikzpicture}
		\end{center}

		Regardless of the formalism, the different approaches must reproduce the same classical predictions. The geometric formulation of GR ultimately leads to the Einstein field equations.
		\newpage
	\subsubsection{A Theory of Gravity}
		General Relativity, to be a theory of gravity, must explain gravitational phenomena.

		For example, Newtonian gravity became a theory of gravity because it explained the results Kepler had found. Kepler described the motion of the planets and discovered that their orbits were elliptical using data analysis, but he did not have an underlying physical theory explaining \textit{why} they followed these trajectories.

		Newton instead derived the observed behavior from first principles. His laws of motion and universal gravitation explained Kepler's laws and provided a general framework for celestial mechanics, and ended up predicting the whole of kinematics and dynamics.

		As such, GR as a theory of gravity must be able to explain observed gravitational phenomena: relativistic corrections to celestial mechanics, gravitational waves, black holes, gravitational collapse, and, on the largest scales, cosmology.

		Perhaps the grandest application of GR is cosmology. Newtonian gravity is fundamentally a local theory applied within a fixed space and time, while GR allows the spacetime itself to be dynamical. This makes it possible to formulate a theory describing the universe as a whole.

		\paragraph{DISTINCTION:
		} For something to be a theory of gravity, what must it possess?

		\begin{enumerate}
			\item \textbf{Dynamical variables}

			The theory must have variables describing the gravitational field in which objects move.

			\begin{align*}
				\text{Newton:} \qquad &\Phi\\
				\text{GR:} \qquad &g_{\mu\nu}
			\end{align*}

			\item \textbf{A law describing how the field affects objects}

			The theory must describe how its dynamical variables determine the motion of objects.

			For Newtonian gravity:
			\begin{equation}
				\vb{a} = -\grad\Phi.
			\end{equation}

			For GR, freely falling objects follow geodesics:
			\begin{equation}
				\dv[2]{x^\sigma}{\lambda} 
				+
				\Gamma^\sigma_{\mu\nu}
				\dv{x^\mu}{\lambda} 
				\dv{x^\nu}{\lambda} 
				=0,
			\end{equation}
			or, using the four-velocity $u^\mu$,
			\begin{equation}
				u^\mu\grad_\mu u^\nu = 0.
			\end{equation}

			\item \textbf{A law describing how objects produce the dynamical variable}

			Finally, the theory must tell us how matter produces the gravitational field.

			For Newtonian gravity:
			\begin{equation}
				\lpc\Phi = 4\pi G\rho.
			\end{equation}

			For GR:
			\begin{equation}
				R_{\mu\nu}
				-\frac{1}{2}g_{\mu\nu}R
				=
				8\pi G T_{\mu\nu}.
			\end{equation}
			Since the two formalisms attempt to explain the same phenomenon, we demand
			that the GR formulation reduces to the Newtonian formulation in weak fields.
		\end{enumerate}
		\par
		\noindent{\textit{\textbf{Remark : }}}
		\par
		Notice that $\lpc \Phi = 4\pi G \rho$ and $
				R_{\mu\nu}
				-\frac{1}{2}g_{\mu\nu}R
				=
				8\pi G T_{\mu\nu}
				$ 
		are inherently similar. They basically describe the same thing, but in different fields of geometry.
		Notice
		\begin{align*} 
			\underbrace{\lpc\Phi}_{\text{some form of }\partial^2} &= \underbrace{4\pi G\rho}_{\text{constant multiplied to matter distribution}}
			\\
			\underbrace{R_{\mu\nu}
			-\frac{1}{2}g_{\mu\nu}R}
			_{\text{some combination of }\partial^2}
																   &=
			\underbrace{8\pi G T_{\mu\nu}}
			_{\text{constant multiplied by matter distribution}}
		\end{align*}

		These three ingredients are the basic structure of a theory of gravity: a dynamical variable, a law governing how it affects matter, and a field equation governing how matter produces it.

		GR is inherently nonlinear. Going from the field to the motion of objects is already more complicated than in Newtonian gravity, but the field equation itself is also nonlinear: the metric determines the curvature, while the curvature determines the metric.

		Solving the Einstein field equations for realistic systems is therefore often extremely difficult. In many cases, exact solutions cannot be obtained, and the equations must be solved numerically. Careers have been made doing this. This is the domain of \textit{numerical relativity}.

		In the next chapters, we will explore how gravity can be explained by general relativity.
